\noindent{\color{stewartblue}\textsf{\textbf{VÍ DỤ 5}}}\quad Một trong những đại lượng được quan tâm trong nhiệt động lực học là độ nén. Nếu một chất cho trước được giữ ở nhiệt độ không đổi, thì thể tích $V$ của nó phụ thuộc vào áp suất $P$. Chúng ta có thể xét tốc độ thay đổi của thể tích theo áp suất---cụ thể là đạo hàm $dV/dP$. Khi $P$ tăng, $V$ giảm, do đó $dV/dP < 0$. \textbf{Độ nén} được định nghĩa bằng cách thêm vào một dấu trừ và chia đạo hàm này cho thể tích $V$:
\begin{equation*}
\text{độ nén đẳng nhiệt} = \beta = -\frac{1}{V} \frac{dV}{dP}
\end{equation*}

Như vậy, $\beta$ đo lường tốc độ giảm thể tích của một chất trên mỗi đơn vị thể tích khi áp suất tác dụng lên nó tăng ở nhiệt độ không đổi.

Chẳng hạn, thể tích $V$ (tính bằng mét khối) của một mẫu không khí ở $25^\circ\text{C}$ được xác định là có liên hệ với áp suất $P$ (tính bằng kilopascal) bởi phương trình
\begin{equation*}
V = \frac{5.3}{P}
\end{equation*}

Tốc độ thay đổi của $V$ theo $P$ khi $P = 50\text{ kPa}$ là
\begin{equation*}
\left. \frac{dV}{dP} \right|_{P=50} = \left. -\frac{5.3}{P^2} \right|_{P=50} = -\frac{5.3}{2500} = -0.00212\text{ m}^3/\text{kPa}
\end{equation*}

Độ nén tại áp suất đó là
\begin{equation*}
\beta = \left. -\frac{1}{V} \frac{dV}{dP} \right|_{P=50} = \frac{0.00212}{\dfrac{5.3}{50}} = 0.02\text{ (m}^3/\text{kPa)}/\text{m}^3 \tag*{{\color{stewartblue}\rule{1.2ex}{1.2ex}}}
\end{equation*}

\vspace{1.8em}

\noindent{\color{stewartblue}\rule{1.2ex}{1.2ex}}\hspace{0.6em}{\sffamily\bfseries\large Sinh học}

\vspace{0.8em}

\noindent{\color{stewartblue}\textsf{\textbf{VÍ DỤ 6}}}\quad Giả sử $n = f(t)$ là số cá thể trong một quần thể động vật hoặc thực vật tại thời điểm $t$. Sự thay đổi kích thước quần thể giữa các thời điểm $t = t_1$ và $t = t_2$ là $\Delta n = f(t_2) - f(t_1)$, và do đó tốc độ tăng trưởng trung bình trong khoảng thời gian $t_1 \leqslant t \leqslant t_2$ là
\begin{equation*}
\text{tốc độ tăng trưởng trung bình} = \frac{\Delta n}{\Delta t} = \frac{f(t_2) - f(t_1)}{t_2 - t_1}
\end{equation*}

\textbf{Tốc độ tăng trưởng tức thời} nhận được từ tốc độ tăng trưởng trung bình này bằng cách cho khoảng thời gian $\Delta t$ tiến dần về $0$:
\begin{equation*}
\text{tốc độ tăng trưởng} = \lim_{\Delta t \to 0} \frac{\Delta n}{\Delta t} = \frac{dn}{dt}
\end{equation*}

Nói một cách chặt chẽ, điều này chưa hoàn toàn chính xác vì đồ thị thực tế của hàm số kích thước quần thể $n = f(t)$ sẽ là một hàm bậc thang gián đoạn bất cứ khi nào có một cá thể sinh ra hoặc chết đi, và do đó không khả vi. Tuy nhiên, đối với một quần thể động vật lớn
